%!TeX program = xelatex
\documentclass[12pt,hyperref,a4paper,UTF8]{ctexart}
\usepackage{UCASReport}
\usepackage{booktabs}
\usepackage{longtable}
\usepackage{caption}
\usepackage{setspace}
\usepackage{float}
\usepackage{seqsplit}
\hypersetup{colorlinks=true,linkcolor=blue,citecolor=blue,urlcolor=blue}
\setlength{\parindent}{2em}
\setlength{\headheight}{15pt}
\setlength{\emergencystretch}{3em}
\setlength{\parskip}{0pt}
\linespread{1.45}
\captionsetup{font=small,labelsep=quad}
\begin{document}
\cover
\thispagestyle{empty}
\clearpage
\pagenumbering{Roman}
\begin{abstract}
随着智慧校园建设不断推进，课堂教学过程逐渐由经验观察向数据辅助分析转变。传统课堂行为记录主要依赖教师人工观察或课后回看视频，存在工作量大、主观性较强、难以形成连续统计等问题。针对上述问题，本文设计并实现了一套基于深度神经网络的课堂学生行为识别与教学互动分析系统。系统以单阶段目标检测网络YOLO26为核心，将课堂中的学生行为划分为书写、阅读、抬头听讲、转头、举手、站立和讨论七类，在公开课堂行为数据基础上完成数据审计、类别重构、场景感知划分、模型训练、对比实验以及图片和视频推理模块开发。\par
数据审计结果表明，所使用的SCBehavior数据实际包含400幅图像和8083个有效目标框。为降低同一课堂场景相邻画面跨集合分布造成的数据泄漏风险，本文利用ResNet18图像特征进行场景相似度分析，将数据划分为281幅训练图像、59幅验证图像和60幅独立测试图像。在相同测试集上，基线YOLO26n@640模型取得0.7711的mAP@0.5和0.5667的mAP@0.5:0.95。通过比较增大模型容量和提高输入分辨率两种改进方案，最终采用YOLO26s@640作为系统模型，其精确率、召回率、mAP@0.5和mAP@0.5:0.95分别达到0.7373、0.7926、0.8221和0.6065。与基线相比，mAP@0.5提高5.10个百分点，阅读、转头和举手等关键类别的AP@0.5分别提高10.92、11.79和12.74个百分点。\par
在工程实现方面，本文使用Python、PyTorch和Ultralytics构建训练及推理流程，开发了支持单幅图像和课堂视频输入的识别模块。系统能够输出带中文行为标签的可视化结果，并生成JSON汇总和CSV时间线，为课堂活动回顾和教学互动趋势分析提供辅助信息。研究同时指出，当前系统输出属于逐帧检测结果，尚未实现跨帧身份跟踪，不能将累计检测次数直接解释为去重后的学生人数或教学质量评分。实验结果表明，该系统能够在有限课堂数据条件下实现较为有效的多类别行为检测，具有一定的课程教学演示价值和后续扩展空间。\par
\par\noindent\textbf{关键词：}智慧校园；课堂行为识别；目标检测；深度学习；YOLO26；教学互动分析
\end{abstract}
\begin{center}\textbf{Abstract}\end{center}
With the development of smart campuses, classroom observation is gradually moving from manual experience toward data-assisted analysis. Traditional observation and retrospective video review are labor-intensive, subjective, and unsuitable for continuous statistics. This paper designs and implements a classroom student behavior recognition and instructional interaction analysis system based on deep neural networks. The system adopts the one-stage YOLO26 detector and recognizes seven behaviors: writing, reading, looking up, turning the head, raising a hand, standing, and discussing. The complete workflow covers dataset auditing, category reconstruction, scene-aware splitting, model training, comparative experiments, and image/video inference.\par
The audited SCBehavior dataset contains 400 images and 8,083 valid bounding boxes. To reduce leakage caused by visually similar neighboring classroom frames appearing in different subsets, ResNet18 embeddings are used for scene similarity analysis, producing 281 training images, 59 validation images, and 60 independent test images. On the same test set, the YOLO26n@640 baseline obtains 0.7711 mAP@0.5 and 0.5667 mAP@0.5:0.95. After comparing model scaling and resolution scaling, YOLO26s@640 is selected as the final model. It achieves precision 0.7373, recall 0.7926, mAP@0.5 0.8221, and mAP@0.5:0.95 0.6065. Its AP@0.5 for reading, head turning, and hand raising improves by 10.92, 11.79, and 12.74 percentage points over the baseline, respectively.\par
The implementation is developed with Python, PyTorch, and Ultralytics. It accepts images and classroom videos, produces visualized results with Chinese labels, and exports JSON summaries and CSV timelines. The outputs are frame-level detections rather than identity-tracked student counts and should therefore be interpreted as auxiliary classroom activity evidence rather than direct teaching-quality scores. The results demonstrate the feasibility of multi-class classroom behavior detection under a limited-data setting and provide a foundation for future temporal modeling and campus deployment.\par
\par\noindent\textbf{Key words:}smart campus; classroom behavior recognition; object detection; deep learning; YOLO26; instructional interaction analysis
\newpage
\tableofcontents
\newpage
\pagenumbering{arabic}
\setcounter{page}{1}
\section{绪论}
\subsection{研究背景}
课堂是高校教学活动最主要的场景之一。学生在课堂中表现出的书写、阅读、抬头听讲、举手和讨论等行为，能够从不同侧面反映教学活动的组织方式和互动状态。以往教师通常依靠现场观察形成总体印象，或者在教学检查中人工回看录像。此类方法能够理解复杂语境，但难以对长时间、多人数课堂进行稳定记录，而且观察结果可能受到观察者经验、注意力和评价标准差异的影响。\par
摄像设备和GPU计算平台在校园中的普及，使计算机视觉辅助课堂分析具备了实现条件。与仅对整幅图像给出单一类别的图像分类不同，课堂画面通常同时包含多名学生，每名学生又可能处于不同状态，因此更适合建模为多目标检测问题：系统不仅需要判断行为类别，还需要给出行为主体在画面中的位置。YOLO系列将目标位置与类别预测整合到统一网络中，具有端到端、推理速度快和部署流程相对简洁等特点，适合课程设计中完成从模型训练到应用验证的完整闭环。\par
\subsection{研究意义}
本课题的意义主要体现在三个方面。第一，在教学辅助层面，识别结果能够帮助教师快速定位举手、讨论等互动片段，减少人工浏览完整视频的时间。第二，在智慧校园实践层面，课题将深度学习、数据治理、模型评估和视频处理结合起来，形成具有具体校园场景的人工智能应用案例。第三，在课程学习层面，项目并非只调用预训练模型完成演示，而是对数据真实性、集合划分、模型误差和方案取舍进行系统验证，有利于理解机器学习实验的可重复性和边界。\par
需要强调的是，课堂行为不等同于学习效果。短时间低头可能是阅读，也可能是记笔记；转头可能表示分心，也可能属于正常讨论。本文将系统定位为课堂活动记录和趋势分析工具，不使用单帧检测结果评价个体学生，更不将其作为自动考核依据。这一定位既符合当前模型能力，也有助于避免技术结论超出证据范围。\par
\subsection{国内外研究现状}
目标检测方法大体可分为两阶段方法和单阶段方法。两阶段方法先生成候选区域，再对候选区域进行分类与边界框回归，通常具有较强的定位能力；单阶段方法直接在特征图上完成密集预测，具有较高的推理效率。Redmon等提出的YOLO将检测建模为从整幅图像到边界框和类别概率的统一回归问题，为实时单阶段检测奠定了基础[1]。Lin等针对密集检测中的前景与背景不平衡问题提出Focal Loss，使模型更加关注难分类样本[3]。残差网络通过学习残差映射缓解深层网络优化困难，已成为视觉特征提取的重要基础[2]。\par
课堂行为研究通常包括三类技术路线：基于人体姿态关键点的方法、基于目标检测的方法以及结合时间信息的视频行为识别方法。姿态方法对举手、站立等具有明显肢体结构的行为较为直观，但在学生密集、遮挡严重、人体尺度较小时容易出现关键点缺失。视频时序方法能够利用前后帧区分瞬时动作，但标注成本和计算量较高。目标检测方法可以直接输出多名学生的空间位置和行为类别，适合在数据规模有限时建立基础系统，因此本文选用该路线。\par
现有课堂行为项目常见问题包括数据规模描述与实际文件不一致、随机划分导致相邻帧泄漏、仅报告总体准确率而缺少逐类别分析，以及只展示检测图片而缺少可复现的测试结果。本文围绕这些问题展开工作：首先审计数据文件和标注；其次采用场景感知方式划分集合；再次设置独立测试集比较模型；最后实现可复用的图片与视频命令行推理模块。\par
\subsection{主要研究内容}
本文主要完成以下工作：一是对SCBehavior数据进行文件级和标注级审计，建立七类课堂行为体系；二是提取图像视觉特征并进行相似度约束的数据划分；三是构建YOLO26n基线，分别从模型容量和输入分辨率角度开展改进实验；四是对总体指标、逐类别指标和错误样本进行分析；五是开发图片与视频推理程序，输出中文标注、行为统计和时间线；六是讨论系统应用边界与进一步改进方向。\par
\section{相关理论与关键技术}
\subsection{卷积神经网络}
卷积神经网络通过局部连接和权值共享提取图像空间特征。对于输入特征图，卷积核在空间位置滑动并计算加权和，随后经过归一化和非线性激活得到新的特征表示。浅层特征通常描述边缘、颜色和纹理，深层特征逐渐形成对人体轮廓、姿态和课堂上下文的语义表达。目标检测网络通过多尺度特征融合，使模型能够同时处理近处大目标与远处小目标。\par
随着网络加深，直接学习复杂映射可能产生梯度传播和优化困难。ResNet引入跨层快捷连接，将目标映射改写为残差函数与恒等映射之和[2]。本文没有直接使用ResNet作为最终检测器，而是使用ImageNet预训练的ResNet18提取全局图像嵌入，用于判断课堂图片之间的场景相似性，从而辅助数据划分。\par
\subsection{目标检测基本任务}
目标检测需要同时解决分类与定位。一个预测结果通常表示为类别、置信度以及边界框坐标。训练过程中，模型通过分类损失学习行为类别，通过边界框回归损失学习目标位置。预测阶段需要根据置信度阈值过滤低可信结果，并处理同一目标产生的重复框。交并比（Intersection over Union，IoU）定义为预测框与真实框交集面积除以并集面积，是衡量定位重合程度的基本指标。\par
传统非极大值抑制按照置信度保留候选框，并移除与高分框高度重叠的重复预测。YOLO26默认提供端到端的一对一检测头，并保留可选的一对多检测路径[4]。考虑到本系统七类行为之间具有互斥倾向，应用层还启用了跨类别重叠抑制，以减少同一学生同时出现多个高度重合行为框的情况。该处理属于面向本任务的后处理约束，不改变训练模型参数。\par
\subsection{YOLO26模型}
YOLO继承了单阶段检测的基本思想：输入图像经过主干网络提取特征，颈部结构融合不同尺度信息，检测头输出类别与位置。YOLO26是Ultralytics在2026年发布的实时视觉模型系列，提供n、s、m、l和x等不同规模，并支持检测、分割、姿态和分类等任务[4]。官方资料将其主要特征概括为端到端推理、轻量化边界框回归以及更新后的训练策略。\par
本文选择YOLO26n和YOLO26s进行实验。n模型参数量较少，适合作为基线；s模型在仍可由RTX 3080高效训练和推理的前提下具有更强特征表达能力。两者均采用COCO预训练权重进行迁移学习，再在课堂行为数据上微调。迁移学习能够利用通用图像中形成的边缘、纹理、人体和物体特征，降低小规模数据从头训练的难度。\par
\subsection{模型评价指标}
精确率P表示模型预测为某类的目标中真实正确的比例，召回率R表示真实目标中被成功检出的比例，可分别写为：\par
\begin{equation}
P=\frac{TP}{TP+FP},\qquad R=\frac{TP}{TP+FN}.
\end{equation}
其中TP、FP和FN分别表示真正例、假正例和假负例。平均精度AP由某一类别的精确率--召回率曲线综合计算，mAP则为各类别AP的平均值。mAP@0.5使用0.5作为IoU阈值，反映较宽松定位要求下的检测能力；mAP@0.5:0.95在0.5至0.95的多个IoU阈值上求平均，对边界框定位质量要求更严格。本文同时报告P、R、mAP@0.5和mAP@0.5:0.95，避免用单一指标概括模型性能。\par
\section{数据集构建与预处理}
\subsection{数据来源与审计}
本文选用GitHub公开项目提供的SCBehavior课堂行为数据作为原始数据来源，并将下载版本固定到提交df45开头的特定版本，以便复现实验。实际审计发现，数据目录包含400幅可用图像和8083个有效目标框，与部分项目说明中出现的1346幅图像不一致。因此本文所有统计和结论均以实际下载文件为准，不沿用未经验证的网页描述。\par
审计过程包括图像可读性检查、标注文件对应关系检查、类别编号合法性检查、边界框坐标范围检查以及空标注统计。经过处理，系统采用书写、阅读、抬头听讲、转头、举手、站立和讨论七个类别。类别选择兼顾数据中已有标注与课堂解释性，没有人为增加缺少训练样本的新类别。\par
\subsection{数据特点}
课堂图像具有目标密集、遮挡明显和行为差异细微等特点。同一画面中可能包含数十名学生，后排学生所占像素较少；阅读、书写和抬头听讲主要依赖头部方向、手部位置及桌面物体关系，视觉边界并不总是清晰；教室座位和拍摄角度相对固定，又容易使模型学习场景背景而非行为本身。此外，各类别样本数量不完全均衡，少数类别测试指标可能因少量样本变化产生较大波动。\par
预处理将原始标注转换为YOLO格式，每行包含类别编号、目标中心点横纵坐标以及宽高，坐标均按图像尺寸归一化。训练时由Ultralytics完成缩放、批处理和常规增强。本文没有通过复制测试图像或人工挑选简单样本提高结果，测试集在模型训练前固定并仅用于最终评估。\par
\subsection{场景感知数据划分}
若直接随机划分课堂图像，同一视频或相邻时刻的画面可能同时进入训练集和测试集。由于背景、座位和学生位置几乎相同，这会使测试结果偏高。为降低该风险，本文使用预训练ResNet18提取每幅图像的特征向量，计算图像之间的余弦距离，并结合场景相似关系完成分组划分。\par
最终训练集、验证集和测试集分别包含281、59和60幅图像，占总数约70.25\%、14.75\%和15.00\%。检查结果显示，在设定的0.01余弦距离阈值内，不存在跨集合的近重复图像对。训练集用于参数学习，验证集用于观察收敛和选择最佳权重，测试集用于报告最终结果。该划分方式不能保证消除所有语义相似场景，但比逐图完全随机划分更符合课堂连续画面的实际特点。\par
\section{系统需求分析与总体设计}
\subsection{功能需求}
系统需要完成四项核心功能。第一，能够加载训练完成的课堂行为检测权重，并自动选择CUDA或CPU设备。第二，能够对单幅图片进行检测，保存带中文标签、置信度和边界框的结果图，同时输出结构化JSON。第三，能够处理常见课堂视频，按设定步长分析视频帧，生成保持原始帧率和总帧数的标注视频。第四，能够输出逐分析帧的CSV时间线和全视频类别累计汇总，方便进一步绘图和查询。\par
\subsection{非功能需求}
系统应具有可重复性、可解释性和容错能力。训练参数、数据划分和模型权重需要保存；输出应同时包含可视化图像与机器可读数据；在没有GPU的设备上能够回退到CPU；输入路径错误或视频无法打开时应给出明确异常。由于课程设计仅要求形成报告，本文没有开发复杂Web前端，而是以清晰的Python命令行脚本验证核心功能，减少与研究目标无关的界面工作。\par
\subsection{系统总体流程}
系统离线阶段依次完成数据审计、集合划分、模型训练、验证集选优和独立测试。在线推理阶段依次完成图像或视频读取、抽帧和尺寸变换、YOLO26s行为检测、置信度筛选与重叠处理、中文标注、逐类计数以及文件保存。该设计将模型训练与应用推理解耦，便于替换权重、调整阈值或扩展输出格式。\par
\begin{figure}[!htbp]
\centering
\includegraphics[width=0.98\textwidth]{figures/06_system_flow.png}
\caption{课堂学生行为识别系统总体流程}
\end{figure}
\subsection{软件与硬件环境}
项目运行于Windows环境，使用Python 3.11.15、PyTorch 2.10.0+cu128和Ultralytics 8.4.91。硬件为NVIDIA GeForce RTX 3080 10GB显存。项目采用独立虚拟环境管理依赖，数据、配置、实验、模型、输出、脚本和文档分别存放，避免将临时文件与正式结果混合。训练脚本和评估脚本均允许通过参数指定模型、图像尺寸、轮次和设备。\par
\section{模型训练与改进方法}
\subsection{基线实验设计}
基线采用YOLO26n预训练权重，输入分辨率为640\$\textbackslash{}times\$640。模型包含约237.62万参数，计算量约5.2 GFLOPs，权重文件约5.12 MB。选择轻量模型作为基线有两个原因：其一，能够快速验证数据和训练管线；其二，后续可明确观察增大模型容量或提高输入尺寸带来的收益与成本。\par
在正式训练前，项目先进行小轮次冒烟测试，检查数据路径、标签读取、CUDA计算、反向传播、验证和权重保存是否正常。冒烟测试通过后再执行完整训练，从而避免长时间运行后才发现配置错误。训练过程中保存每轮损失和验证指标，并以验证表现最佳的权重进行独立测试。\par
\subsection{基线误差分析}
基线在测试集上取得0.6944的精确率、0.7796的召回率、0.7711的mAP@0.5和0.5667的mAP@0.5:0.95。从逐类别结果看，讨论类别AP@0.5达到0.9329，站立为0.8643，书写和抬头听讲分别为0.8273和0.8484；阅读和转头仅为0.5543和0.6044，是主要薄弱类别。\par
错误样本分析表明，核心问题并非只由小目标造成，而是阅读、抬头听讲与转头之间存在较强语义混淆。这些行为在静态图像中的差别集中于头部角度、视线和手部局部关系，轻量模型的特征表达能力可能不足。基于这一观察，本文设计两条改进路线：一是将模型由n扩大到s，增强特征表达；二是保持n模型但将输入尺寸提高到960，以保留更多小目标细节。\par
\subsection{改进方案一：增大模型容量}
第一种方案采用YOLO26s@640。其参数量约946.79万，计算量20.5 GFLOPs，权重约19.36 MB。与基线相比，输入尺寸和数据划分保持不变，主要变量为模型规模，因此能够较清楚地分析容量变化。该模型训练耗时285.76 s，峰值显存约4.78 GiB，RTX 3080能够稳定完成训练。\par
\subsection{改进方案二：提高输入分辨率}
第二种方案采用YOLO26n@960，仍使用约237.62万参数，但计算量随输入尺寸提高到11.7 GFLOPs。更高分辨率理论上有利于后排学生、手部和头部细节，但会增加显存和训练时间。实测该方案训练耗时532.83 s，约为基线的1.91倍，峰值显存约3.10 GiB。\par
\subsection{最终模型选择原则}
模型选择不只比较单一最高mAP，还需要考虑关键类别表现、推理效率、训练成本和应用标准化。YOLO26n@960的总体mAP略高于YOLO26s@640，但差值很小；YOLO26s在阅读、抬头听讲、转头和举手类别上表现更好，而且保持常用的640输入尺寸。综合考虑后，本文选择YOLO26s@640作为最终模型，并将其权重固定保存用于图片和视频推理。\par
\section{实验结果与分析}
\subsection{总体性能对比}
三种模型的独立测试结果如下表所示。\par
\begin{table}[!htbp]
\centering
\small
\resizebox{\textwidth}{!}{%
\begin{tabular}{lccccccc}
\toprule
模型 & 输入尺寸 & 参数量/M & P & R & mAP@0.5 & mAP@0.5:0.95 & 推理时间/ms \\
\midrule
YOLO26n & 640 & 2.376 & 0.6944 & 0.7796 & 0.7711 & 0.5667 & 6.3143 \\
YOLO26s & 640 & 9.468 & 0.7373 & 0.7926 & 0.8221 & 0.6065 & 6.9635 \\
YOLO26n & 960 & 2.376 & 0.7798 & 0.7669 & 0.8257 & 0.6083 & 7.0858 \\
\bottomrule
\end{tabular}}
\caption{不同实验方案综合性能对比}
\label{tab:model-comparison}
\end{table}
最终模型相对基线的mAP@0.5提高0.0510，即5.10个百分点，相对增幅约6.61\%；mAP@0.5:0.95提高0.0399，即3.99个百分点，相对增幅约7.04\%。精确率和召回率也分别提高4.29和1.30个百分点，说明改进不仅增加检出数量，也减少了一部分错误预测。\par
\begin{figure}[!htbp]
\centering
\includegraphics[width=0.92\textwidth]{figures/01_overall_metrics.png}
\caption{不同模型在独立测试集上的检测性能对比}
\end{figure}
\subsection{逐类别性能分析}
最终YOLO26s@640在书写、阅读、抬头听讲、转头、举手、站立和讨论类别上的AP@0.5分别为0.8883、0.6635、0.9069、0.7223、0.8934、0.7512和0.9288。相对基线，书写提高6.09个百分点，阅读提高10.92个百分点，抬头听讲提高5.85个百分点，转头提高11.79个百分点，举手提高12.74个百分点。这验证了增大模型容量对若干细粒度行为具有明显帮助。\par
同时，站立类别下降11.31个百分点，讨论下降0.41个百分点。站立类别在测试集中的样本规模有限，少量漏检会造成较大指标波动；更大模型也可能对局部姿态和背景产生不同偏好。该现象说明总体mAP提升不代表每一类都提升，后续应补充站立行为数据、进行类别均衡采样并检查难例标注。\par
\begin{figure}[!htbp]
\centering
\includegraphics[width=0.92\textwidth]{figures/02_per_class_ap50.png}
\caption{不同模型各行为类别AP@0.5对比}
\end{figure}
\begin{figure}[!htbp]
\centering
\includegraphics[width=0.92\textwidth]{figures/03_final_vs_baseline_gain.png}
\caption{最终模型相对基线的逐类别性能变化}
\end{figure}
\subsection{精度与效率分析}
YOLO26s参数量约为基线的3.98倍，但在本机测试中的模型推理时间仅由6.31 ms增加到6.96 ms。该时间不包含完整的视频解码、绘图与文件写入，因此不能直接等价为端到端系统帧率。YOLO26n@960虽然参数量不变，但训练时间显著增加，说明高分辨率带来的计算成本主要发生在更大的特征图上。对于本课程设计使用的RTX 3080，三种模型均可运行；若部署到低功耗设备，则应重新基准测试并考虑模型导出、量化或使用轻量基线。\par
\begin{figure}[!htbp]
\centering
\includegraphics[width=0.82\textwidth]{figures/04_accuracy_speed_tradeoff.png}
\caption{模型精度—速度权衡}
\end{figure}
\subsection{训练收敛与混淆分析}
训练曲线显示三种方案的验证集mAP总体随训练推进上升，并在后期趋于稳定，说明迁移学习能够在有限课堂数据上有效收敛。不同轮次之间仍存在波动，这与数据规模较小、行为边界模糊和类别不均衡有关。独立测试集的归一化混淆矩阵进一步显示，阅读、转头和抬头听讲之间仍是主要混淆来源。仅依靠单帧外观难以完全判断视线和动作目的，后续引入时序信息具有合理性。\par
\begin{figure}[!htbp]
\centering
\includegraphics[width=0.96\textwidth]{figures/05_training_curves.png}
\caption{三种方案验证集mAP训练曲线}
\end{figure}
\begin{figure}[!htbp]
\centering
\includegraphics[width=0.8\textwidth]{figures/07_final_confusion_matrix_normalized.png}
\caption{最终模型归一化混淆矩阵}
\end{figure}
\subsection{结果可信度讨论}
本文通过固定数据版本、保存划分清单、分离验证集与测试集、对三种模型使用相同测试集以及保留原始评估文件提高实验可信度。但400幅图像仍属于小规模数据，测试集只有60幅图像，指标不能代表所有教室、机位、光照和学生群体。模型在当前数据分布上的有效性不能直接推广到其他学校，实际部署前应采集目标教室数据并进行隐私合规评估。\par
\section{系统实现与功能验证}
\subsection{推理模块结构}
核心推理类BehaviorDetector负责加载最终权重、选择设备、执行预测、转换检测结果和绘制中文标签。图片脚本接收输入路径、输出目录、置信度阈值和设备参数，保存标注图与JSON。视频脚本读取视频元数据，逐帧处理并写出标注视频，同时在CSV中记录分析帧编号、时间和七类检测数量，在JSON中记录总体配置和累计统计。\par
为保证视频播放连续，未被选作推理的间隔帧仍按原图写入输出视频，而不是重复上一分析帧。该设计避免使用抽帧加速时出现画面冻结。系统还保留原视频帧率和总帧数，使输出视频时长与输入基本一致。\par
\subsection{图片识别流程}
图片进入模型后被缩放到设定输入尺寸，网络输出边界框、类别和置信度。程序将坐标映射回原图，在目标框上绘制中文类别名称与置信度，并统计七类行为数量。JSON记录输入文件、模型、阈值、图像尺寸和逐目标结果，便于后续程序读取。\par
测试图片能够正确显示中文字体，没有出现方框或乱码。识别结果同时包含多名学生及不同类别，说明模型与标注绘制链路正常。定性结果只能展示系统是否能够工作，性能结论仍以完整测试集的量化指标为准。\par
\begin{figure}[!htbp]
\centering
\includegraphics[width=0.9\textwidth]{figures/09_system_demo.jpg}
\caption{系统单幅图片中文标注效果}
\end{figure}
\subsection{视频识别流程}
视频模块使用OpenCV获取分辨率、帧率和总帧数，根据frame\_stride参数决定分析间隔。被分析帧通过YOLO26s获得检测结果并绘制标签，所有帧按原顺序写入MP4。测试使用30帧、10 fps的示例视频，设置步长2，共分析15帧；输出视频仍包含30帧且能够正常读取，证明抽帧分析与视频保存逻辑正确。\par
时间线CSV每一行对应一个被分析帧，可用于观察不同时间点行为数量变化。汇总JSON中的累计次数是各分析帧检测数量之和，同一学生跨多帧出现会被重复计数。若要统计独立学生人数，需要加入多目标跟踪和身份关联模块，当前版本不具备该能力。\par
\subsection{功能测试结果}
环境检查确认PyTorch能够识别RTX 3080并使用CUDA训练和推理；模型权重能够正常加载；图片输出、JSON输出、视频输出、CSV时间线和汇总JSON均成功生成；输出视频的帧数和帧率经程序重新读取验证。最终模型文件还计算了SHA-256摘要，用于确认后续报告和演示使用的是同一权重版本。\par
\section{系统应用价值、局限性与改进方向}
\subsection{校园应用价值}
系统适合用于课堂录像的辅助检索。例如，教师可根据举手或讨论行为的时间线快速定位互动片段，回顾问题设计和课堂节奏；教学研究人员可在取得授权后，对匿名化课堂样本进行群体活动趋势分析；人工智能课程可将其作为数据集处理、模型训练、评估和部署的一体化案例。系统的价值在于提供可复核的视觉线索，而不是替代教师作出教育判断。\par
\subsection{当前局限性}
第一，数据规模较小，来源场景有限，模型对新教室和新机位的泛化能力尚未验证。第二，阅读、书写、抬头听讲与转头在静态画面中存在语义重叠，标签本身也可能受标注者判断影响。第三，系统未使用时序特征，不能识别需要连续动作才能定义的复杂行为。第四，视频累计统计没有身份去重。第五，遮挡、光照变化、后排小目标和摄像机抖动仍可能导致漏检。第六，课堂视频涉及个人信息，真实部署需要遵守数据最小化、授权、访问控制和保存期限等要求。\par
\subsection{后续改进方向}
后续可从数据、模型和系统三个层面改进。数据层面可采集不同教室、课程、光照和机位的数据，补充站立等少数类并开展交叉场景测试；模型层面可加入更适合小目标的检测层、注意力机制、知识蒸馏或类别均衡策略，也可将目标检测与姿态估计结合；视频层面可引入ByteTrack等多目标跟踪方法，并使用时序网络或Transformer融合连续帧；应用层面可开发只呈现群体统计的可视化界面，加入数据脱敏与权限管理，避免存储不必要的个人身份信息。\par
\section{总结与展望}
本文完成了一套面向校园课堂环境的学生行为识别系统。项目从公开数据出发，对实际文件和标注进行审计，确认400幅图像和8083个目标框；建立七类行为体系，并利用ResNet18特征完成场景感知划分；训练YOLO26n基线并针对语义混淆问题设计YOLO26s@640和YOLO26n@960两组改进实验；最终选择YOLO26s@640作为系统模型。\par
独立测试结果表明，最终模型mAP@0.5达到0.8221，mAP@0.5:0.95达到0.6065，相对基线分别提高5.10和3.99个百分点。阅读、转头和举手等关键类别改善明显，但站立类别出现下降，说明小规模和不均衡数据下仍存在类别波动。工程实现支持图片和视频识别，能够生成中文标注图、标注视频、JSON汇总与CSV时间线，完成了从数据处理、神经网络训练到应用输出的课程设计闭环。\par
总体而言，深度学习目标检测可以为课堂活动记录提供有效辅助，但模型输出需要结合具体情境谨慎解释。未来工作应扩大跨场景数据、融合姿态与时间信息、加入多目标跟踪，并将隐私保护纳入系统设计。只有在技术指标、教育解释和合规边界三方面共同完善，课堂行为识别才适合从课程原型进一步走向实际校园应用。\par
\clearpage
\begin{thebibliography}{99}
\bibitem{ref1} REDMON J, DIVVALA S, GIRSHICK R, et al. You Only Look Once: Unified, Real-Time Object Detection[C]//Proceedings of the IEEE Conference on Computer Vision and Pattern Recognition. 2016: 779-788.
\bibitem{ref2} HE K, ZHANG X, REN S, et al. Deep Residual Learning for Image Recognition[C]//Proceedings of the IEEE Conference on Computer Vision and Pattern Recognition. 2016: 770-778.
\bibitem{ref3} LIN T Y, GOYAL P, GIRSHICK R, et al. Focal Loss for Dense Object Detection[C]//Proceedings of the IEEE International Conference on Computer Vision. 2017: 2980-2988.
\bibitem{ref4} JOCHER G, QIU J, LIU M, et al. Ultralytics YOLO26: Unified Real-Time End-to-End Vision Models[EB/OL]. arXiv:2606.03748, 2026.
\bibitem{ref5} EVERINGHAM M, ESLAMI S M A, VAN GOOL L, et al. The Pascal Visual Object Classes Challenge: A Retrospective[J]. International Journal of Computer Vision, 2015, 111(1): 98-136.
\bibitem{ref6} LIN T Y, MAIRE M, BELONGIE S, et al. Microsoft COCO: Common Objects in Context[C]//European Conference on Computer Vision. 2014: 740-755.
\bibitem{ref7} PASZKE A, GROSS S, MASSA F, et al. PyTorch: An Imperative Style, High-Performance Deep Learning Library[C]//Advances in Neural Information Processing Systems. 2019, 32.
\bibitem{ref8} BRADSKI G. The OpenCV Library[J]. Dr. Dobb's Journal of Software Tools, 2000.
\bibitem{ref9} BEWLEY A, GE Z, OTT L, et al. Simple Online and Realtime Tracking[C]//2016 IEEE International Conference on Image Processing. 2016: 3464-3468.
\bibitem{ref10} ZHANG Y, SUN P, JIANG Y, et al. ByteTrack: Multi-Object Tracking by Associating Every Detection Box[C]//European Conference on Computer Vision. 2022: 1-21.
\end{thebibliography}
\appendix
\section{主要程序文件}
训练程序为\path{scripts/train_yolo.py}，测试集评估程序为\path{scripts/evaluate_yolo.py}，实验比较程序为\path{scripts/compare_experiments.py}，图片推理程序为\path{scripts/predict_image.py}，视频推理程序为\path{scripts/predict_video.py}，核心推理实现为\path{src/campus_behavior/inference.py}。项目环境和使用命令将在正式排版稿附录中进一步整理。\par
\section{实验可重复性信息}
数据源版本固定到提交df45开头的版本；数据划分清单位于\path{data/splits/scbehavior_manifest.csv}；最终模型为\path{models/final/classroom_behavior_yolo26s.pt}；其SHA-256为\texttt{\seqsplit{F4FF75828F4160DF2E544D10F17F5EC30E8E17A93AB4FE6415B3C4F107F36BF8}}。所有正式指标来自独立测试集，不使用系统演示图片代替量化评估。\par
\end{document}
